\documentclass[11pt]{article}

\usepackage[utf8]{inputenc}
\usepackage[T1]{fontenc}
\usepackage{amsmath,amssymb,amsthm}
\usepackage{booktabs}
\usepackage{geometry}
\usepackage{microtype}
\usepackage[hidelinks]{hyperref}
\geometry{a4paper,margin=1in}

\newtheorem{theorem}{Theorem}
\newtheorem{lemma}{Lemma}
\newtheorem{corollary}{Corollary}
\newtheorem{proposition}{Proposition}
\theoremstyle{remark}
\newtheorem{remark}{Remark}

\newcommand{\Dsw}{D_{\mathrm{sw}}}
\newcommand{\Sthree}{S_3}

\title{Two faces of one optimum: wobble error-resistance and minimal
tRNA decoding coincide in the genetic code}
\author{Jonathan Reich\\ \normalsize Independent Researcher}
\date{\today}

\begin{document}
\maketitle

\begin{abstract}
The structure of the standard genetic code has been read as optimal in two
separate literatures: it is highly error-resistant (synonymous codons cluster
under single mutations), and it is decoded by a remarkably small set of tRNAs.
We prove these are the same optimum. The wobble-fibre constancy that uniquely
maximises the third-position silent-edge count---the dominant component of
single-substitution error-resistance---also uniquely minimises the number of
tRNAs needed to decode the code under four-way (superwobble) reading: the
minimiser of the decoding number is exactly the maximiser of silent edges. This
coincidence is specific to \emph{superwobble}---the documented four-way decoding
in which a single unmodified tRNA reads an entire four-codon family, the
mechanism by which organelles and reduced genomes translate with minimal tRNA
sets \cite{rogalski}. Far from a technical caveat, this specificity is the
theory's empirical engine: it predicts that the joint optimum is approached as
genomes reduce toward superwobble decoding and departed from as tRNA repertoires
expand. We give exact statements and proofs, validate the decoding model (it
reproduces the classical 31-tRNA Crick minimum and the 22-tRNA mitochondrial
set), and confirm the prediction across 390 genomes spanning the three domains:
decoding redundancy above the minimum scales with genome size (pooled
$p\approx2\times10^{-18}$; significant within all three domains; and
$p\approx6\times10^{-10}$ under phylogenetic generalised least squares), with the
most reduced genomes nearest the minimum. The result bears on a long-standing
dispute---whether the code's error-resistance reflects direct selection for
robustness \cite{freeland} or is a non-adaptive by-product \cite{massey,koonin}:
if one structural property delivers error-resistance and decoding economy
together, neither trait need be the other's selective cause, and the
adaptive/neutral dichotomy is, for this component of robustness, ill-posed.
\end{abstract}

\paragraph{Significance.} Why the genetic code is so resistant to mutation is a
classic question, answered either by selection for error-tolerance or by appeal
to neutral evolution. We prove a theorem that makes the dispute moot for a major
component of robustness: the feature that renders the code maximally
error-resistant under third-position mutation---grouping codons into four-codon
families---is \emph{identical} to the feature that lets it be decoded by the
fewest possible tRNAs. Error-resistance and translational economy are one
optimum, not two. The link is realised through superwobble decoding, and we
confirm its population-level signature across 390 genomes spanning all domains of
life.

\section{Introduction}

The near-universal genetic code is redundant in a strikingly non-random way.
Two features have each been called optimal. First, \emph{error-resistance}:
synonymous codons differ predominantly at the third (wobble) position, so most
single-nucleotide substitutions are silent or conservative, and the code beats
almost all random codes on error-load measures \cite{freeland,haig}. Second,
\emph{decoding economy}: although there are $61$ sense codons, organisms read
them with far fewer tRNA species---$31$ at the Crick minimum \cite{crick66},
and as few as $22$ in animal mitochondria---by exploiting wobble pairing at the
anticodon \cite{agris,grosjean}. These two optimalities have been studied
independently, with error-resistance framed in terms of amino-acid assignment
and decoding economy in terms of tRNA modification and gene content.

Here we show they are not independent: they are two expressions of the same
combinatorial property. A companion result \cite{wobble} proves that the
third-position silent-edge count $\Sthree$ is uniquely maximised by
\emph{wobble-fibre constancy}---making each two-base codon prefix wholly
synonymous (a four-codon family box). We prove that wobble-fibre constancy also
uniquely \emph{minimises} the number of tRNAs required to decode the code under
superwobble reading \cite{rogalski}. The argmin of the decoding number equals the
argmax of $\Sthree$ (Theorem~\ref{thm:dual}). The result is exact, elementary, and
machine-verified, and it has a sharp biological corollary: because the
coincidence holds under four-way superwobble (Remark~\ref{rmk:crick})---the
regime in which organelles and reduced genomes actually decode whole family boxes
with single, unmodified tRNAs \cite{rogalski}---the lineages in which robustness
and economy are \emph{jointly} optimised are exactly the reduced,
superwobble-using genomes, while expanded cytoplasmic repertoires, decoding
through modified wobble, sit further from the optimum. We confirm the resulting
prediction---that decoding redundancy above the minimum grows with genome
size---in a 390-genome, phylogenetically controlled comparative survey.

The conceptual payoff speaks to one of the oldest questions about the code.
Freeland and Hurst's demonstration that the standard code beats almost all
alternatives on error load---``one in a million'' \cite{freeland,haig}---is
widely read as evidence of selection \emph{for} mutational robustness; a
neutralist tradition counters that robustness emerged as a by-product of code
expansion and tRNA coevolution \cite{massey,koonin}. The dual optimality shows
the framing is under-determined. A single structural property---large family
boxes---yields error-resistance and a minimal decoding apparatus simultaneously,
so any force that enlarges family boxes (selection for tRNA economy, genome
reduction, biased substitution) delivers error-resistance for free, and
conversely. Optimality in one currency \emph{is} optimality in the other; neither
need be the proximate target of selection.

\section{Results}

\subsection{The silent-edge functional and the decoding number}

Let $\mathcal{C}=\{A,C,G,U\}^3$ be the $64$ codons; a \emph{code} is a map
$f\colon\mathcal{C}\to\Sigma$ onto the $20$ amino acids and \texttt{Stop}. Two
codons are adjacent if they differ at one position; an edge is \emph{silent} if
its endpoints share $f$. Let $\Sthree(f)$ count silent edges whose substitution
is at the third position. Writing $\pi(c)$ for the two-base prefix of $c$, the
$16$ prefixes partition $\mathcal{C}$ into \emph{boxes} of four codons, and two
codons share a prefix iff they are third-position-adjacent; hence the
third-position adjacency graph is $16$ disjoint copies of $K_4$. With
$n_{\pi,a}=\#\{c:\pi(c)=\pi,f(c)=a\}$,
\begin{equation}\label{eq:S3}
  \Sthree(f)=\sum_{\pi}\sum_{a}\binom{n_{\pi,a}}{2}.
\end{equation}

A tRNA charges one amino acid and reads codons that differ only at the third
position, i.e.\ within one box. Under \emph{superwobble} a single modified
anticodon can read any subset of a box's third positions \cite{agris,grosjean};
thus one tRNA decodes each non-empty $(\pi,a)$ group ($a\neq\texttt{Stop}$), and
the minimal decoding number is
\begin{equation}\label{eq:Dsw}
  \Dsw(f)=\sum_{\pi} m_\pi,\qquad m_\pi=\#\{a\neq\texttt{Stop}:n_{\pi,a}>0\}.
\end{equation}
Under strict Crick wobble a single anticodon reads only certain pairs, and a
four-codon family costs two tRNAs; we return to this in
Remark~\ref{rmk:crick}.

\subsection{Dual optimality}

\begin{lemma}[Per-box monotonicity]\label{lem:box}
Within a box, merging two amino-acid groups of sizes $a,b$ raises that box's
contribution to $\Sthree$ by $ab>0$ and lowers its contribution to $\Dsw$ by
$1$. Hence, per box, $\Sthree$ is maximised iff $\Dsw$ is minimised, at the
monochromatic partition.
\end{lemma}
\begin{proof}
By \eqref{eq:S3} the box contributes $\sum_i\binom{k_i}{2}$;
$\binom{a+b}{2}-\binom{a}{2}-\binom{b}{2}=ab$. By \eqref{eq:Dsw} it contributes
its number of groups, which falls by $1$ on a merge.
\end{proof}

\begin{theorem}[Dual optimality]\label{thm:dual}
For every code, $\Dsw(f)\ge16$ and $\Sthree(f)\le96$, and the following are
equivalent: \emph{(i)} $\Dsw(f)=16$; \emph{(ii)} $\Sthree(f)=96$; \emph{(iii)}
$f$ is wobble-fibre constant. Consequently
$\arg\min_f\Dsw(f)=\arg\max_f\Sthree(f)$: the minimal-tRNA codes are exactly the
maximally wobble-robust codes.
\end{theorem}
\begin{proof}
Each of the $16$ boxes contributes $\ge1$ to $\Dsw$ and $\le\binom{4}{2}=6$ to
$\Sthree$, both extremes reached together iff the box is monochromatic
(Lemma~\ref{lem:box}). Sum over boxes.
\end{proof}

\begin{remark}[Superwobble specificity is a prediction, not a fragility]\label{rmk:crick}
Under Crick wobble a monochromatic box costs two tRNAs (one each for $\{C,U\}$
and $\{A,G\}$), as does a $2{+}2$ pyrimidine/purine split; the decoding number
no longer distinguishes family boxes from such splits and the coupling of
Theorem~\ref{thm:dual} fails. The dual optimality is thus a property of the
\emph{superwobble} regime---not an idealisation but the documented four-way
decoding by which unmodified $U_{34}$ tRNAs read entire family boxes, enabling
translation with reduced tRNA sets in organelles and minimal genomes
\cite{rogalski,agris}. Its regime-specificity is therefore falsifiable rather
than fragile: it predicts that codes jointly optimising robustness and economy
are those of superwobble-using lineages---paradigmatically animal mitochondria,
combining family-box structure with the minimal $22$-tRNA decoding set---whereas
genomes decoding through modified wobble carry larger, more redundant
repertoires.
\end{remark}

\begin{proposition}[Constrained co-monotonicity]\label{prop:constr}
Among codes with a fixed multiset of block sizes, both $\Sthree$ and $-\Dsw$ are
Schur-convex in the per-box group-size profile: any majorising transfer does not
decrease $\Sthree$ and does not increase $\Dsw$. The block structures favoured by
error-resistance are exactly those favoured by decoding economy; codes rich in
four-codon families, like the standard code's eight family boxes, are jointly
near-optimal.
\end{proposition}
\begin{proof}[Proof sketch]
$\Sthree$ sums the Schur-convex $(k_i)\mapsto\sum_i\binom{k_i}{2}$ over boxes;
$\Dsw$ counts non-empty parts, non-increasing under majorising transfers.
\end{proof}

\subsection{The standard code and the full code space}

Computing $\Dsw$ from \eqref{eq:Dsw} reproduces the classical Crick minimum of
$31$ tRNAs for the standard code, and $22$ for vertebrate mitochondria under
superwobble---validating the decoding model against known biology. The standard
code is not wobble-fibre constant (it has split boxes such as $AAR=$Lys,
$AAY=$Asn); it attains neither global extremum but sits at a constrained
optimum, as encoding $20$ amino acids and \texttt{Stop} forbids the degenerate
joint optimum of Theorem~\ref{thm:dual} (Proposition~\ref{prop:constr}). Across
$3000$ random codes the ``soft'' coupling holds, $\mathrm{corr}(\Sthree,\Dsw)
\approx-0.55$, and the standard code sits in the jointly favourable corner (high
$\Sthree=64$, low $\Dsw=23$) relative to the random mean ($\Sthree\approx15$,
$\Dsw\approx57$).

\subsection{Decoding redundancy scales with genome size}

Theorem~\ref{thm:dual} and Remark~\ref{rmk:crick} predict that genomes lie above
the theoretical minimum decoding set by a \emph{redundancy factor}
$R=(\text{distinct tRNA anticodons})/\Dsw^{\min}$ that shrinks toward $1$ in
reduced, superwobble-leaning genomes. We tested this in a sample of $390$ genomes
spanning bacteria, archaea and eukaryotes (GtRNAdb tRNA gene sets; NCBI assembly
sizes; Methods). Mean decoding repertoire was $37.4$ (bacteria), $42.5$
(archaea) and $44.2$ (eukaryotes) against the standard-code superwobble minimum
of $23$. Regressing $R$ on $\log_{10}$ genome size over the $310$ genomes with
resolved sizes gave a positive slope ($0.152$ per decade,
$p\approx1.9\times10^{-18}$, Spearman $\rho=+0.47$). Because genome size and domain
are correlated, we controlled the confound three ways. \emph{(i)} Adding domain
as a covariate leaves the size effect strong ($0.21$ per decade,
$p\approx8.7\times10^{-15}$), so the relationship is not merely the
prokaryote/eukaryote contrast. \emph{(ii)} \emph{Within} each domain separately
the slope is positive and significant in all three: bacteria $0.38$ per decade
($p\approx4\times10^{-7}$, $\rho=+0.42$, $n=102$, $0.3$--$10$\,Mb), eukaryotes
$0.18$ ($p\approx9\times10^{-7}$, $\rho=+0.73$, $n=106$, $2.5$\,Mb--$4.8$\,Gb)
and archaea $0.23$ ($p\approx9\times10^{-3}$, $n=102$, $0.5$--$5$\,Mb). Reduced
genomes lie nearest the minimum---the obligate intracellular endosymbionts
\emph{Buchnera} ($0.64$\,Mb) and \emph{Riesia} ($0.58$\,Mb), \emph{Rickettsia},
\emph{Wolbachia} and the reduced methanogens fall to repertoires of $30$--$35$,
approaching the minimum of $23$---as predicted by Remark~\ref{rmk:crick}.
\emph{(iii)} Finally, phylogenetic generalised least squares on the
NCBI-taxonomy tree of the $309$ taxa gives a positive, strongly significant slope
under \emph{every} branch-length model: $0.26$ per decade under Brownian motion
with unit edges ($t=6.4$, $p\approx6\times10^{-10}$), $0.29$ under Grafen branch
lengths \cite{grafen89} ($p\approx8\times10^{-9}$), and $0.25$ at the
maximum-likelihood Pagel's $\lambda=0.83$ \cite{pagel99}
($p\approx6\times10^{-9}$). The fitted $\lambda$ reports strong phylogenetic
signal in decoding redundancy (likelihood-ratio test against $\lambda=0$:
$p\approx5\times10^{-25}$), yet the size effect is undiminished and, if anything,
steeper than the naive OLS slope ($0.15$). The relationship is therefore robust
to phylogenetic-tree specification, not an artefact of non-independence.

\section{Discussion}

The standard genetic code's error-resistance and its tRNA decoding economy are
the same optimum. The wobble-fibre constancy that maximises third-position
silent edges is exactly the property that minimises the superwobble decoding
set (Theorem~\ref{thm:dual}); away from the optimum the two track one another
(Proposition~\ref{prop:constr}; $\mathrm{corr}\approx-0.55$); and empirically,
decoding redundancy grows with genome size as the superwobble corollary predicts.

The main consequence is for the long debate over whether the code's
error-minimisation is adaptive \cite{freeland,haig} or a non-adaptive by-product
\cite{massey,koonin}. The dual optimality shows that, for the third-position
component of error-resistance, the question is under-determined at the structural
level: large family boxes confer error-resistance and a minimal decoding
apparatus inseparably, so any force that enlarges them---selection for tRNA
economy, genome reduction, biased substitution---produces error-resistance as an
automatic correlate, and conversely. Robustness here is plausibly a free rider on
decoding economy, or the reverse, rather than an independent selective target.
More broadly, the result is a concrete instance of \emph{structural pleiotropy}:
when two purported ``optimalities'' of a biological code are governed by a single
combinatorial parameter, demonstrating optimality in one currency is not evidence
of selection in that currency, and apparent adaptive fine-tuning can be the
shadow of one constraint. The same caution applies wherever code- or
network-level optimality is read adaptively without first checking whether the
optimised quantities are structurally independent.

A natural objection is that superwobble is an idealisation: most cytoplasmic
tRNAs carry $U_{34}$ modifications that \emph{restrict} wobble, so that a
four-codon family is typically read by more than one tRNA. This is precisely why
the coupling has empirical content. Superwobble is the documented decoding mode
of organelles and genome-reduced bacteria \cite{rogalski,agris}, the very systems
that translate with minimal tRNA sets; modified-wobble cytoplasmic systems decode
the same families with several tRNAs and therefore sit \emph{above} the joint
optimum. The theory thus predicts not a uniform fact but a gradient---exact dual
optimality in the superwobble regime, growing redundancy as genomes expand and
shift to modified wobble---and that gradient is what the $390$-genome survey
detects: mitochondria attain the $22$-tRNA minimum, reduced endosymbionts
approach it, and repertoires climb monotonically with genome size. A property
that held under every decoding model would forbid this signature;
superwobble-specificity is what makes the dual optimality testable, and the test
passes.

\paragraph{Limitations.} The decoding model counts a minimal tRNA cover under
fixed wobble rule sets, abstracting away modification chemistry, tRNA abundance
and decoding efficiency, which shape realised repertoires; the genome-size
analysis uses distinct anticodons as a repertoire proxy rather than
expression-weighted decoding. The phylogenetic control uses the
NCBI-taxonomy topology; we show the slope is robust across branch-length models
(unit, Grafen, and a maximum-likelihood Pagel's $\lambda=0.83$) and that decoding
redundancy carries strong phylogenetic signal, but the topology is taxonomic
rather than inferred from molecular sequences, so a time-calibrated species tree
remains a refinement---one unlikely to overturn an effect significant at
$p\approx6\times10^{-9}$ under every model tested. The
silent-edge functional is a deliberately simple, unweighted measure of
error-resistance; physicochemically weighted measures \cite{haig,freeland} would
test how far the structural pleiotropy extends. Finally, the theorem concerns the
third-position component $\Sthree$, which dominates but does not exhaust
single-substitution error-resistance.

\section{Methods}

\paragraph{Functional and decoding number.} $\Sthree$ and $\Dsw$ are computed
from \eqref{eq:S3}--\eqref{eq:Dsw} on the standard and NCBI variant codes. The
minimal decoding number per $(\pi,a)$ group is the minimum number of wobble rule
sets covering the group's third-position bases; under the superwobble rule set
(which includes a four-way reader) this is $1$ per group, giving \eqref{eq:Dsw}.
The Crick rule set ($G_{34}\!:\!\{C,U\}$, $U_{34}\!:\!\{A,G\}$,
$I_{34}\!:\!\{A,C,U\}$, $C_{34}\!:\!\{G\}$, $A_{34}\!:\!\{U\}$) reproduces the
$31$-tRNA minimum.

\paragraph{Comparative data.} tRNA gene sets were obtained from GtRNAdb by
scraping the genome index, downloading each genome's tRNAscan-SE results and
counting distinct anticodons (mapped to cognate codons); the decoding repertoire
is the number of distinct anticodons. Genome sizes were obtained from NCBI
Assembly (\texttt{esearch}/\texttt{esummary}, total assembly length). The
redundancy factor is repertoire divided by the standard-code superwobble minimum
($23$). Regression of $R$ on $\log_{10}$ genome size was by ordinary least
squares (pooled, with domain as a covariate, and within each domain) and by
phylogenetic generalised least squares under a Brownian-motion variance-covariance
matrix derived from the NCBI-taxonomy tree of the sampled taxa (genomes resolved
to taxids via NCBI; lineages assembled into a tree; strains sharing a taxid
averaged). To check robustness to branch-length specification, the PGLS was run
under unit edges, Grafen \cite{grafen89} branch lengths and a maximum-likelihood
Pagel's $\lambda$ transform \cite{pagel99}; the fitted $\lambda$ and a
likelihood-ratio test against $\lambda=0$ (no phylogenetic signal) are reported.
Spearman correlation is reported alongside. All code and data fetchers
are scripted (Python; \texttt{numpy}/\texttt{statsmodels}); the engine
reproduces $\Sthree=64$, $S=69$ and the $31$/$22$-tRNA minima as regression
tests.

\begin{thebibliography}{99}

\bibitem{wobble} [Author], ``Sharp combinatorial constraints on wobble
robustness and coding architecture in the genetic code,'' under revision,
\emph{Mathematical Biosciences} (2026).

\bibitem{crick66} F.~H.~C.~Crick, ``Codon--anticodon pairing: the wobble
hypothesis,'' \emph{Journal of Molecular Biology} \textbf{19} (1966) 548--555.

\bibitem{grafen89} A.~Grafen, ``The phylogenetic regression,'' \emph{Philosophical
Transactions of the Royal Society B} \textbf{326} (1989) 119--157.

\bibitem{pagel99} M.~Pagel, ``Inferring the historical patterns of biological
evolution,'' \emph{Nature} \textbf{401} (1999) 877--884.

\bibitem{agris} P.~F.~Agris, F.~A.~P.~Vendeix and W.~D.~Graham, ``tRNA's wobble
decoding of the genome: 40 years of modification,'' \emph{Journal of Molecular
Biology} \textbf{366} (2007) 1--13.

\bibitem{grosjean} H.~Grosjean, V.~de Cr\'ecy-Lagard and C.~Marck, ``Deciphering
synonymous codons in the three domains of life: co-evolution with specific tRNA
modification enzymes,'' \emph{FEBS Letters} \textbf{584} (2010) 252--264.

\bibitem{rogalski} M.~Rogalski, D.~Karcher and R.~Bock, ``Superwobbling
facilitates translation with reduced tRNA sets,'' \emph{Nature Structural \&
Molecular Biology} \textbf{15} (2008) 192--198.

\bibitem{freeland} S.~J.~Freeland and L.~D.~Hurst, ``The genetic code is one in
a million,'' \emph{Journal of Molecular Evolution} \textbf{47} (1998) 238--248.

\bibitem{haig} D.~Haig and L.~D.~Hurst, ``A quantitative measure of error
minimization in the genetic code,'' \emph{Journal of Molecular Evolution}
\textbf{33} (1991) 412--417.

\bibitem{massey} S.~E.~Massey, ``Genetic code evolution reveals the neutral
emergence of mutational robustness, and information as an evolutionary
constraint,'' \emph{Life} \textbf{5} (2015) 1301--1332.

\bibitem{koonin} E.~V.~Koonin and A.~S.~Novozhilov, ``Origin and evolution of
the genetic code: the universal enigma,'' \emph{IUBMB Life} \textbf{61} (2009)
99--111.

\end{thebibliography}

\end{document}
